\documentclass[11pt]{article}
\usepackage[margin=25mm]{geometry}
\usepackage{fontspec}
\setmainfont{Latin Modern Roman}
\usepackage{amsmath,amssymb,mathtools}
\usepackage{unicode-math}
\setmathfont{Latin Modern Math}
\usepackage{xurl,hyperref}
\hypersetup{colorlinks=true,urlcolor=blue}
\setcounter{secnumdepth}{0}
\setlength{\parindent}{0pt}
\setlength{\parskip}{5pt}
\setlength{\emergencystretch}{3em}
\newcommand{\sref}[2]{\hyperlink{source:#1:#2}{[S#2]}}
\newcommand{\cataloguescope}[1]{\paragraph{Recorded target.}#1}
\newcommand{\R}{\mathbb R}
\newcommand{\C}{\mathbb C}
\newcommand{\N}{\mathbb N}
\newcommand{\Z}{\mathbb Z}
\newcommand{\Q}{\mathbb Q}
\newcommand{\F}{\mathbb F}
\newcommand{\eps}{\varepsilon}
\begin{document}
\section{Asymptotic Metrics for Higher-Rank Monopoles}\label{statement}

\subsection{Definitions and mathematical statement}

\paragraph{Shared notation.}$\mathbb N=\{1,2,\ldots\}$; the standard number fields have their usual meanings. All additional conventions are stated locally.


\paragraph{Definitions and conventions.} A finite-energy monopole on $\mathbb R^3$ consists of a connection $A$ and Higgs field $\Phi$ for the gauge group, in the Bogomolny/monopole setting of the source. Framed moduli divide by gauge transformations tending to the identity at infinity. Magnetic charge is determined by the asymptotic Higgs field: one positive integer $k$ for $SU(2)$, or a vector of $N-1$ charges for maximally broken $SU(N)$. The moduli metric is the $L^2$ inner product on infinitesimal field variations perpendicular to gauge directions. Maximal symmetry breaking means the limiting Higgs value is regular, so its centralizer is a maximal torus. The source compactification adds boundaries describing clusters of lower-charge monopoles separating to infinity. Quasi-isometric metric comparison on an asymptotic region means two metrics bound one another by positive constants there, with the covering and cluster-coordinate conventions retained.

\paragraph{Source scope and qualifications.} The target includes arbitrary permitted magnetic-charge vectors and the compactification's cluster strata, rather than a single separated-charge example. The source calls this work in progress; local metric approximations alone do not yield the requested uniform quasi-isometric control across all relevant regions.

\paragraph{Catalogue mathematical target.}
\paragraph{Catalogue wording retained for traceability.}
For $\operatorname{SU}(N)$-monopoles with maximal symmetry breaking, establish the quasi-isometric asymptotics of the monopole metric on the analogous compactification and determine whether they agree qualitatively with the $\operatorname{SU}(2)$ case.

\paragraph{Saved literature context.} The supplied catalogue's assessment dated 2026-08-21 records: Asymptotic metrics for maximally broken SU(N) monopoles are known, but the proposed compactification and the claimed global quasi-isometric control for arbitrary charges remain unresolved; the record itself describes this as work in progress. The assessment is retained as dated provenance; it is not a new comprehensive literature or open-status certification.
\subsection{Short English statement}
Control the higher-rank monopole metric near every compactification boundary and compare it with the SU(2) cluster asymptotics.
\subsection{Research attempt}
\paragraph{Model, date and investigated scope.}
OpenAI GPT-6 Astra (configured model \texttt{gpt-6-astra}), ultra reasoning; generated 11 October 2026. Proves a relative-operator metric comparison criterion and a degenerating-coordinate counterexample to an insufficient error norm.

\paragraph{Sources and literature check.}
The complete supplied definition, including source qualifications and the dated literature assessment, was reviewed. The calculation below is an elementary restricted attempt motivated by [S2]; it does not claim novelty or a fresh literature-status determination.

\paragraph{Context and investigated scope.} The complete supplied definition, including source qualifications and the dated literature assessment, was reviewed. The calculation below is an elementary restricted attempt motivated by [S2]; it does not claim novelty or a fresh literature-status determination.

\paragraph{The exact relative-error condition for metric comparison.} Let $g_0$ be a positive metric on an asymptotic region and let $g=g_0+h$. Define the self-adjoint endomorphism $A=g_0^{-1/2}hg_0^{-1/2}$ in a local frame. If $\|A\|\leq\epsilon<1$ uniformly, then for every tangent vector $v$
\[
 (1-\epsilon)g_0(v,v)\leq g(v,v)\leq(1+\epsilon)g_0(v,v).
\]
Indeed put $w=g_0^{1/2}v$ and bound $|w^TAw|\leq\epsilon|w|^2$. A relative estimate $\|A\|\leq C/R$ gives a uniform two-sided comparison for $R\geq2C$.

\paragraph{Small coordinate errors are not enough.} Take $g_0=\operatorname{diag}(R^{-4},1)$ and $h=\operatorname{diag}(R^{-2},0)$ for $R\to\infty$. Every entry of $h$ tends to zero, but $g(e_1,e_1)/g_0(e_1,e_1)=1+R^2\to\infty$. A collapsing phase direction can therefore amplify an apparently small asymptotic error. Controlling only the separated-position block would miss it.

\paragraph{Monopole gap.} This is a concrete criterion for the proposed quasi-isometric comparison, not an estimate for the actual higher-rank metric. One must show relative bounds simultaneously in position, phase and cluster directions, uniformly as several compactification faces meet and for each allowed magnetic-charge vector. Existing pointwise leading terms cannot be promoted by the coordinate-smallness argument just disproved. Establishing those uniform relative errors and a compatible compactification is the first unresolved step before comparing qualitatively with the $SU(2)$ ends.

\paragraph{Verification and adversarial audit.}
Checked congruence transformation, epsilon<1 positivity, R≥2C constants and the explicit anisotropic counterexample. Model self-audit only; no independent human proof checking. Prompts A and B were applied to this canonical record and its qualified source scope. Checked congruence transformation, epsilon<1 positivity, R≥2C constants and the explicit anisotropic counterexample. Retained unresolved status and zero original-target progress; Uniform relative metric estimates across all higher-rank cluster strata and charges are missing.

Coordinator checked the relative metric quadratic-form bound and failure of coordinatewise smallness in a collapsing direction. This additional check was model-assisted, not independent human verification.

\paragraph{Tools, process and contributions.}
Python 3 was used for file preparation and structural checks; no numerical experiment or theorem prover was used. Contribution prompt A was followed by prompt B. GPT-6 Astra with ultra reasoning drafted the research and performed its stated self-audit; any additional coordinator contribution is identified in the verification and audit. A GPT-6-based Codex coordinator assembled the manuscript and checked repository integration; its exact variant and reasoning configuration were not exposed. No human mathematical contribution, independent human proof check or formal proof certification is supplied. Complete derivations and the consulted source references are included here; private reasoning transcripts are not published.

\paragraph{Final status and first remaining gap.}
\textbf{UNRESOLVED}. Uniform relative metric estimates across all higher-rank cluster strata and charges are missing.

\section{Completion Estimate}
\noindent\textbf{COMPLETION: 0\%}
The calculation is an auxiliary or restricted attempt and establishes no newly verified portion of the full source target. Uniform relative metric estimates across all higher-rank cluster strata and charges are missing.

\subsection{Sources}
\begin{itemize}
\item[\hypertarget{source:statement:1}{S1}] ULAM.AI and the UnsolvedMath Contributors, UnsolvedMath, catalogue entry. Catalogue curation under CC BY 4.0; definitions and English exposition added locally.
\url{https://huggingface.co/datasets/ulamai/UnsolvedMath}.
\item[\hypertarget{source:statement:2}{S2}] Analysis and Geometric Singularities (2012), Oberwolfach Reports, DOI 10.4171/owr/2012/25, pp. 1487--1562; official report source and mathematical context.
\url{https://doi.org/10.4171/owr/2012/25}.
\item[\hypertarget{source:statement:3}{S3}] Official EMS Press PDF of the same Oberwolfach report; PDF page 21, printed page 1507.
\url{https://ems.press/content/serial-article-files/46396}.
\item[\hypertarget{source:statement:4}{S4}] Bielawski, Asymptotic metrics for SU(N)-monopoles with maximal symmetry breaking; reference in the saved dated catalogue literature review.
\url{https://arxiv.org/abs/hep-th/9801092}.
\item[\hypertarget{source:statement:5}{S5}] Kottke--Singer, Partial compactification of monopoles and metric asymptotics; reference in the saved dated catalogue literature review.
\url{https://arxiv.org/abs/1512.02979}.
\end{itemize}
\label{end:statement}
\end{document}
